\section{Logarithmic compression on demand}\label{sec:compression}

Calibration requests the sizes of particular minors as it chooses pivots.
The choices must also be consistent with minors that never occur in a
pivot calculation, since such minors can enter a circuit. This is the
extension requirement in Theorem~\ref{thm:compression}. Write \(\ell_u\)
for the base-two logarithm of the magnitude of a nonzero minor. We seek
short rational answers to individual requests that extend to a vector
of the same linear type as a small perturbation of \(\ell\).

The procedure learns a set \(J\) of minor indices. Between restarts it
uses the logarithms at these indices to set common scales for all its
answers. A failure to represent a requested logarithm at those scales
enlarges \(J\). The bound on the number of restarts comes from the fact
that products of minors are polynomials of bounded degree in the same
matrix entries.

Throughout, an integer is \emph{short} if its bit length is polynomial
in \(d,K\). Fix
\[
 D=(d+1)^6,\qquad
 \eta^{-1}=\text{the least power of two at least }4D\cdot2^K .
\]
We first work with fixed proxies \(v_u\) satisfying
\(|v_u-\ell_u|\le\eta\) and access to the signs of integer affine
forms in these proxies. The final subsection constructs proxies for
which the required sign tests can be implemented in the arithmetic
model of Section~\ref{sec:arithmetic}.

\subsection{A bound on independent logarithms}

Call a set \(I\) of minor indices \emph{dissociated} if the subset sums
of \((v_u)_{u\in I}\) are pairwise at distance greater than two.
Equivalently, every nonzero coefficient vector in
\(\{-1,0,1\}^I\) gives a sum of absolute value greater than two.
This is a version with tolerance of the usual additive notion
\cite{LY}. If \(I\) is maximal, every proxy is within two of a signed
subset sum on \(I\). Indeed, for an index outside \(I\), any relation
preventing its addition must have coefficient \(+1\) or \(-1\) on the
new index.

\begin{lemma}\label{lem:dissociated}
Every dissociated set has size less than \(D\).
\end{lemma}
\begin{proof}
Suppose there is a dissociated set of size \(D\). The logarithms of
the magnitudes of its \(2^D\) subset products of signed minors differ by more
than \(2-D\eta>1\). In increasing order of magnitude each product
therefore has magnitude greater than the sum of all preceding
magnitudes. It follows that all \(2^{2^D}\) subset sums of these
products are distinct at the input matrix. Two unequal subsets
differ at a largest product, whose contribution survives cancellation.

These sums are consequently distinct integer polynomials in the
at most \(d^2\) matrix entries. Each has degree at most \(dD\) and
coefficients of absolute value at most \(2^D(d!)^D\), by the
determinant expansion. Counting possible coefficients at each
monomial gives
\[
 2^{2^D}\le
 \big(1+2^{D+1}(d!)^D\big)^{(dD+1)^{d^2}}.
\]
Two logarithms yield
\(D\le d^2\log(dD+1)+\log(2+D(1+d^2))
\le(7d^2+8)\log(d+1)\), contradicting \(D=(d+1)^6\).
\end{proof}

Goemans's bound on geometrically decreasing sequences of small integer
combinations uses a related counting argument
\cite[Lemma~3.5 in the preprint]{Newton}, attributed there via
Radzik~\cite{Rad}. Here the integer combinations are products of
minors expanded in monomials. Counting subset sums brings the
number of monomials into the dissociation bound only logarithmically.

\subsection{Extending partial answers}

Let \(S\) contain the indices answered since the last restart. To
certify answers on \(S\), we shall enforce agreement of linear type
at a larger coefficient bound \(M\) on \(S\) alone. Lemma~\ref{lem:extend}
shows why this suffices. A maximal dissociated set expresses all the
missing logarithms, to within a rounded offset, in fewer than \(D\)
variables. Projecting out these variables determines \(M\).

\begin{lemma}[Extension]\label{lem:extend}
Set
\[
 A=2^K(D+4)/\eta,\qquad M=(D+1)A^{D+1}.
\]
Suppose on a subset \(S\) of the indices we have
\begin{equation}\label{eq:answered}
 f_S\sim_M g_S,\qquad \|g_S-v_S\|_\infty\le\eta ,
\end{equation}
where \(f_S\) has a short common denominator and short numerators.
Then \(f_S,g_S\) extend to vectors \(f,g\) satisfying
\eqref{eq:compress} on all the minor indices. The vector \(f\) again
has a short common denominator and short numerators.
\end{lemma}
\begin{proof}
Choose a maximal dissociated set \(I\). For every \(u\notin S\),
write \(v_u\) as a signed subset sum on \(I\) plus an offset in
\([-2,2]\). Round the offset to the \(\eta\)-grid, and let \(h_u(x)\)
be the resulting affine form in \(x\in\RR^I\). Set
\(g_u=h_u(v_I)\). On \(S\) and its complement alike,
\(|g_u-\ell_u|\le 2\eta\le2^{-K}\).

Now impose the sign conditions
\[
  \big(y_S,\,(h_u(x))_{u\notin S}\big)\sim_{2^K}g .
\]
Here \(y_S\) is a parameter and \(x\) is an independent variable,
even when \(I\) meets \(S\). Write each sign as a strict inequality or
as two weak inequalities in the zero case. After clearing the grid
denominator, each row has integer coefficients of total absolute
value at most \(A\). For \(y_S=g_S\) we have the solution \(x=v_I\).
Since \(|I|<D\) by Lemma~\ref{lem:dissociated} and \(f_S\sim_M g_S\),
Lemma~\ref{lem:projection} gives a solution also at \(y_S=f_S\).

If \(U\) bounds the common denominator and the absolute numerators
of \(f_S\), the rows after substitution and clearing that denominator
have coefficient norm at most \(AU\). The second assertion of
Lemma~\ref{lem:projection} supplies \(x\) with a common denominator
and absolute numerators at most \((AU+1)^{D+1}\). Set \(f_u=h_u(x)\)
for \(u\notin S\). The product of the two common denominators and
\(1/\eta\) works for the whole vector. Both it and the absolute
numerators are bounded by \((D+2)(AU+1)^{D+1}U/\eta\).
\end{proof}

\subsection{Extracting common scales}

The maximal set \(I\) in Lemma~\ref{lem:extend} serves only to prove
existence. The procedure maintains its own dissociated set \(J\),
starting with the empty set. To decide whether it can answer a
request, it needs a test that succeeds on every proxy within two of
a signed subset sum on \(J\). Enumerating those sums would cost
exponential time. Instead we expand \(v_J\) at a sequence of common
scales. At each scale the vector of digits is integral and the
remaining error is a small fraction of the scale. This is the
simultaneous approximation step of Frank--Tardos rounding~\cite{FT}.

\begin{lemma}\label{lem:peel}
Let \(w\in\RR^m\), and fix integers \(\epsilon^{-1}\ge2\) and
\(H\ge0\). Put \(C=2^{m(m+1)}(2/\epsilon)^{m+1}\).
Given the signs of integer-coefficient affine forms on \(w\), one
can find integer vectors \(z_1,\ldots,z_p\) and positive scales
\(\lambda_1,\ldots,\lambda_p\), represented by rational linear forms
on \(w\), such that
\[
 \begin{aligned}
 &p\le m,\qquad \|z_j\|_\infty\le C,\qquad \lambda_j>H/C,\\
 &\left\|w-\sum_{i=1}^j\lambda_i z_i\right\|_\infty\le\epsilon\lambda_j
       \quad (1\le j\le p),\qquad
   \left\|w-\sum_{i=1}^p\lambda_i z_i\right\|_\infty\le H .
 \end{aligned}
\]
Counting each sign query as one operation, the bit time and the
lengths of all queries and outputs are polynomial in \(m\) and the
bit lengths of \(\epsilon^{-1},H\).
\end{lemma}
\begin{proof}
Start with remainder \(w'=w\). If \(X=\|w'\|_\infty>H\), use binary
search to find a rational vector \(u\) within \(\epsilon/(2C)\) of
\(w'/X\) in each coordinate. Consider the lattice of vectors
\((qu-z,q/C)\) for integral \(q,z\), of determinant \(1/C\).
Lattice reduction~\cite{LLL}, as in~\cite[Theorem~(5.3.19)]{GLS},
produces in polynomial bit time a nonzero vector of Euclidean norm
at most \(2^{m/2}C^{-1/(m+1)}\le\epsilon/2\). Such a vector must
have \(q\ne0\), since a nonzero integer vector has length at least
one. Choose its sign to give \(1\le q\le C\). Then
\[
 \|q w'/X-z\|_\infty\le\epsilon ,\qquad \|z\|_\infty\le q.
\]
For the second bound use integrality and \(\epsilon<1\).

Take \(\lambda=X/q>H/C\) and replace \(w'\) by \(w'-\lambda z\).
The new remainder has norm at most \(\epsilon\lambda\). At a zero
coordinate the digit is zero. At a coordinate attaining \(X\) in
absolute value the digit equals \(+q\) or \(-q\), so that coordinate
now vanishes as well. Thus at most \(m\) steps are needed to reach
residual norm at most \(H\). For \(m=0\) the expansion is empty.

Each comparison for the binary search can be made after multiplying
by \(X>0\). At every stage \(X\) is a signed coordinate of the
remainder, chosen by comparisons, so both the remainders and the
scales are rational linear forms in the original \(w\). A step
multiplies their common denominator by at most \(q\) and increases
the maximum absolute numerator by at most a factor \(2C\). Clearing
these denominators in the comparisons, including those on the
\(\epsilon/(2C)\)-grid, gives sign queries of the asserted length.
\end{proof}

For a homogeneous form, a nonzero integer digit contributes at least
one whole scale. Choosing the residual small compared with that scale
lets the first nonzero digit determine the sign. We record the usual
Frank--Tardos consequence used to round offsets in
Section~\ref{sec:contractions} and, through that computation, to
solve games in Section~\ref{sec:games}.

\begin{lemma}\label{lem:ft}
Given \(w\in\RR^m\) and an integer \(M\ge1\), one can find an integer
vector preserving the sign on \(w\) of every homogeneous linear form
with integer coefficients of \(1\)-norm at most \(M\). The procedure
is strongly polynomial with the bit length of \(M\) counted as
additional input size, and returns the entries in bits.
\end{lemma}
\begin{proof}
Apply Lemma~\ref{lem:peel} with \(\epsilon=1/(2M)\) and \(H=0\)
(cf.\ \cite[Theorem~(6.2.13)]{GLS}). For a form whose first nonzero
digit occurs on \(z_j\), the remaining contribution has magnitude
at most \(M\epsilon\lambda_j<\lambda_j\). A form vanishing on every
digit also vanishes on \(w\), since the final remainder is zero.
Consequently the integer vector
\(\sum_{j=1}^p(2CM+1)^{p-j}z_j\) has the required signs. Its first
nonzero digit dominates the geometric sum of all later digits,
whose evaluations have magnitude at most \(CM\) each. The short
linear forms queried in Lemma~\ref{lem:peel} can be evaluated as in
Section~\ref{sec:arithmetic}.
\end{proof}

For affine forms, append a coordinate equal to one before applying
Lemma~\ref{lem:ft}. Its rounded value is positive by sign agreement.
Dividing the rounded vector by that value preserves the homogeneous
signs and restores the last coordinate to one.

\subsection{Answering a request}

Start each sequence of requests by applying Lemma~\ref{lem:peel} to
\(v_J\). Use \(M\) from Lemma~\ref{lem:extend} and set
\[
 m=|J|,\qquad \epsilon^{-1}=100(m+1)M,\qquad H=100MC ,
\]
where \(C\) is determined by \(m,\epsilon\) in Lemma~\ref{lem:peel}.
All extracted scales \(\lambda_j\) exceed \(100M\). They remain fixed
throughout the sequence.

For a request at \(u\), look for integer digits \(b_{u1},\ldots,b_{up}\)
of magnitude at most \(mC\) satisfying
\begin{equation}\label{eq:request}
 \begin{aligned}
 \left|v_u-\sum_{i=1}^j\lambda_i b_{ui}\right|
       &\le\frac{\lambda_j}{10M} && (1\le j\le p),\\
 \left|v_u-\sum_{i=1}^p\lambda_i b_{ui}\right|
       &\le mH+2 .
 \end{aligned}
\end{equation}
Once the previous digits are fixed, the permitted interval for
\(b_{uj}\) has length \(1/(5M)<1\). Binary search therefore finds the
unique digit or detects failure. Check the last residual bound even if
there are no digits.

\paragraph{A failed request enlarges the dissociated set.}
If \(v_u\) lies within two of a signed subset sum on \(J\), take
the same signed sum of the digit vectors of \(v_J\) at each scale.
This gives digits of magnitude at most \(mC\) with residuals at most
\(m\epsilon\lambda_j+2\le\lambda_j/(10M)\) at position \(j\)
and at most \(mH+2\) at the end. Thus failure implies that \(u\)
lies outside \(J\) and that \(J\cup\{u\}\) is dissociated.
In that case retain the enlarged set, discard the expansion and
previous answers, and declare a restart. Each restart increases
\(|J|\) by one, so Lemma~\ref{lem:dissociated} bounds the total
number by \(D-1\).

\paragraph{A successful request gives a short answer.}
Round the final residual down to the \(\eta\)-grid by binary search,
obtaining \(e_u\) with \(|e_u|\le mH+2\) and error at most \(\eta\).
For \(B=10M(mH+3)\), return \(f_u\) and define its comparison value
\(g_u\) by
\[
 f_u=\sum_{j=1}^p B^{p+1-j}b_{uj}+e_u,\qquad
 g_u=\sum_{j=1}^p\lambda_j b_{uj}+e_u .
\]
Replacing the true scales by powers of \(B\) makes \(f_u\) rational
with common denominator \(1/\eta\) and absolute numerator at most
\(B^{p+2}/\eta\). Retaining the true scales in \(g_u\) gives
\(|g_u-v_u|\le\eta\).

To verify \(f_S\sim_M g_S\), write an integer affine form of norm at
most \(M\) as \(a_0+\sum_{u\in S}a_u x_u\). At position \(j\), its
combined digit is the integer \(\sum_{u\in S}a_u b_{uj}\). If all
combined digits vanish, the two evaluations equal
\(a_0+\sum_{u\in S}a_u e_u\). Otherwise let \(j\) be the first
nonzero combined digit. For \(g_S\), the magnitude contributed
after position \(j\), including the constant term, is bounded by
\[
 M\max\{\lambda_j/(10M)+\eta,1\}<\lambda_j
\]
by \eqref{eq:request} and the rounding error. For \(f_S\), each later
combined digit and the final affine term in \(e\) have magnitude at
most \(B/10\). Their geometric sum is smaller than \(B^{p+1-j}\).
The digit at \(j\) therefore determines both signs. The answers made
in this sequence satisfy \eqref{eq:answered}, and
Lemma~\ref{lem:extend} extends them to all minor indices.

Only \(J\) and its expansion are needed to answer the next request.
The comparisons in \eqref{eq:request} involve at most \(m+1\)
proxies. The product of the denominators \(q\) in Lemma~\ref{lem:peel}
clears coefficients in the scales, and a further factor \(10M/\eta\)
clears the residual and grid comparisons. All these integers
are short since \(p\le m<D\).
Thus initialization and each request require polynomial bit time with
sign access. A uniform integer bound \(F\) for Theorem~\ref{thm:compression}
is available at initialization by using \(U=B^{p+2}/\eta\) in
Lemma~\ref{lem:extend} and taking the largest bound over
\(0\le p\le m<D\).

\subsection{Implementing the sign tests}

A sign test on exact logarithms amounts to comparing a product of
powers of minors with a power of two. An exponent can have short bit
length and still produce very large rational intermediates. We use
rounded logarithms whenever their magnitudes are small. A large
logarithm, when it occurs, certifies that the rational input itself
has sufficient bit length for the power comparison.

Before defining the proxies, fix a short power of two \(T\) such
that every integer sign query used above has coefficient \(1\)-norm
at most \(T\eta\). This bound is uniform over choices of the proxies.
For example, the common denominators and absolute numerators for the
scale coefficients are bounded by \((2C)^m\), and the other
coefficients in the queries were bounded above in terms of \(d,K\).
These estimates allow \(T\) to be fixed for the worst case \(m<D\).

\paragraph{Determining the proxy.}
For each minor needed in a comparison, compute its absolute value
and invert it if necessary to obtain \(z_0=2^{|\ell_u|}\ge1\).
Compute powers \(2^k\) by squaring for \(k=1,2,4,\ldots,T\), stopping
at the first \(k\) for which \(z_0<2^k\). If no test succeeds then
\(|\ell_u|\ge T\) and we take \(v_u=\ell_u\), represented by the
minor itself. Otherwise binary search using powers of two supplies
the integer part of \(\log z_0\). The scalar powers computed in these
tests have exponents at most \(\max\{1,2|\ell_u|\}\). On rational
data this is polynomial in the input length by the determinant bound
of Section~\ref{sec:arithmetic}. The operation count is polynomial
in \(d,K\), since \(T\) is short.

After removing that integer power of two, the remaining factor
\(z\) lies in \([1,2)\). Compute the first \(\log(1/\eta)+6\) terms of
\[
 \ln z=\sum_{j\ge1}\frac{(1-1/z)^j}{j}
\]
and divide by the corresponding partial sum at two. Both series
tails admit a geometric bound, and the partial sum at two is at
least \(1/2\), so the quotient approximates \(\log z\) to error at
most \(\eta/4\). Use binary search to round to the nearest point of
the \(\eta\)-grid, breaking ties deterministically. Restoring the
integer part and sign produces explicit bits for a proxy within
\(\eta\) of \(\ell_u\). The fractional-part computation uses
polynomial-size formulas in \(z\) and short constants, and needs
only comparisons to perform the rounding.

\paragraph{Evaluating an affine form.}
If every participating minor has \(|\ell_u|<T\), evaluate the query
on the explicit proxy bits. Otherwise substitute the rounded proxies
and multiply by \(1/\eta\). The result is an integer constant plus
an integer combination of true logarithms. The constant and each
coefficient have magnitude at most \(2T^2\), by the norm bound
\(T\eta\) on the original form. Exponentiation to base two reduces
the sign test to comparing a product of integer powers of two and
of absolute minor values with one.

Binary powering gives a polynomial operation count in \(d,K\).
For rational data in this branch, some \(|\ell_u|\ge T\), so the
determinant bound puts \(T\) itself within a polynomial in total
input length. Each powering here could therefore be expanded into a
formula with polynomially many leaves in that length, proving the
required bound on intermediate log heights in both arithmetic
conventions of Section~\ref{sec:arithmetic}. Each query uses the
input matrix and short discrete data afresh. This proves
Theorem~\ref{thm:compression} and finishes the proof of
Theorem~\ref{thm:cal}.

\begin{remark}
The dissociation bound uses only the number of input variables, the
degree, and the coefficient \(1\)-norm of the minor polynomials.
The same compression procedure applies to any family of nonzero
values of integer polynomials in at most \(d^2\) scalar inputs whose
degrees are at most \(d\) and coefficient \(1\)-norms at most \(d!\).
It suffices that the indices are short and that each evaluation has
a deterministic strongly polynomial implementation in the chosen
arithmetic convention.
\end{remark}
